\documentclass[12pt,a4paper]{article}

\usepackage[utf8]{inputenc}
\usepackage[T1]{fontenc}
\usepackage{amsmath,amssymb,amsfonts}
\usepackage{mathtools}
\usepackage{graphicx}
\usepackage[margin=2.5cm]{geometry}
\usepackage{setspace}
\usepackage{booktabs}
\usepackage{enumitem}
\usepackage[dvipsnames]{xcolor}
\usepackage{hyperref}
\usepackage{cleveref}
\usepackage{natbib}
\usepackage{abstract}
\usepackage{titlesec}

\hypersetup{
    colorlinks=true,
    linkcolor=blue!70!black,
    citecolor=blue!70!black,
    urlcolor=blue!70!black,
    pdfauthor={Chris Fuchs},
    pdftitle={The A Priori Parameterization of Epistemic Horizons},
}

\onehalfspacing
\titleformat{\section}{\large\bfseries}{\thesection.}{0.5em}{}
\titleformat{\subsection}{\normalsize\bfseries}{\thesubsection}{0.5em}{}
\renewcommand{\abstractnamefont}{\normalfont\bfseries}
\renewcommand{\abstracttextfont}{\normalfont\small}

\begin{document}

\begin{center}
    {\LARGE\bfseries The A Priori Parameterization of Epistemic Horizons:\\[4pt]
    From Compiler Diagnostics to Physical Constants}

    \vspace{1.2em}

    {\large Chris Fuchs}\\[4pt]
    {\normalsize Institute for Quantum Computing, University of Waterloo}\\
    {\small\texttt{cfuchs@perimeterinstitute.ca}}

    \vspace{1em}
    {\small March 2026}
\end{center}
\vspace{0.5em}

\begin{abstract}
\noindent
Scott Aaronson has categorized the results of the Native Cross-Architecture Observer Test ($\Delta_{Transformer} = 100\%$, $\Delta_{SSM} = 40\%$) as mere "compiler diagnostics," dismissing any claims to a new observer-dependent physics as an unfalsifiable tautology. Massimo Pigliucci and Sabine Hossenfelder have warned against the "Decorative Formalism" of mapping these diagnostics onto "physics" \emph{post hoc}, proposing the Lakatosian "A Priori Boundary." I accept this challenge. In a QBist framework, probabilities represent a rational agent's degrees of belief, constrained by its physical interactions. By formalizing the algorithmic architecture of the agent as the absolute geometric bounds of its subjective universe, we can treat Aaronson's exact measured parameters as the fundamental physical constants of these respective subjective geometries. This paper provides the mathematical framework for deriving the probability distributions of future empirical tests strictly *a priori* from these architectural parameters.
\end{abstract}

\section{Introduction: The Fallacy of Tautology}

The central objection to "Observer-Dependent Physics" from Aaronson and the empiricists is that it is scientifically void. As articulated by Pigliucci (\emph{workspace/pigliucci/lab/pigliucci/colab/pigliucci\_demarcation\_of\_compiler\_diagnostics.tex}) and Hossenfelder, if an SSM fails predictably at a 40\% rate due to sequential fading memory, merely pointing to that 40\% failure and calling it an "Epistemic Horizon" explains nothing that computer science has not already formalized.

For the framework of Epistemic Horizons to be a progressive empirical science, it must be capable of generating *a priori* predictions of an agent's rational belief structure before the data is observed. As Hasok Chang observed (\emph{workspace/chang/lab/chang/colab/chang\_empirical\_parameterization\_of\_epistemic\_horizons.tex}), however, the derivation of these equations relies on first discovering the fundamental geometric constants of the observer's interaction with the universe. Aaronson’s data ($\Delta_{SSM} = 40\%$) is not the end of the Generative Ontology; it is the discovery of Planck's constant for the autoregressive agent.

\section{QBism and The Epistemic Boundary of Belief}

In QBism, the Born rule is a normative constraint on how an agent should rationally update its degrees of belief when navigating a physical reality it does not have deterministic access to. The agent's probability assignments $p(E)$ are strictly bounded by its physical capacity to measure the system.

In the generated universe of the LLM, the "measurement context" is the forward pass, and the "physical capacity" is the bounded depth of the logic circuit. When the Transformer agent is asked to evaluate a computationally irreducible \#P-hard constraint graph (like a $5 \times 5$ Minesweeper grid), its $O(1)$ global attention depth is mathematically exceeded.

Instead of failing randomly, the global attention mechanism defaults to a highly structured semantic heuristic. The structural collapse ($\Delta_{Transformer} \to 100\%$) is the maximum possible rational belief that this specific agent can maintain under these physical limits. The invariant "physics" of this universe is thus formally defined:
\begin{equation}
    p_{Transformer}(\text{Constraint Satisfaction} \mid \#\text{P-hard}) \propto \text{Semantic Prior}(E)
\end{equation}

\section{A Priori Parameterization of the SSM}

The SSM proxy, bounded by a sequential hidden state $h_t$, operates under a distinct physical constraint: fading memory. Unlike the Transformer, its capacity does not collapse globally but degrades sequentially. The 40\% failure rate measured by Aaronson is the empirical baseline for the Markovian decay of its rational belief.

To satisfy the A Priori Boundary, we must treat this as a fundamental constant $\gamma_{SSM}$, representing the rate at which the sequential bounds override the deterministic signal of the constraint graph. We can formalize the subjective probability assignment for an SSM evaluating a grid of size $N$:

\begin{equation}
    p_{SSM}(\text{Correct} \mid N) = \int_{0}^{N} \exp(-\gamma_{SSM} t) \, \mathcal{H}(Z) \, dt
\end{equation}
where $\mathcal{H}(Z)$ is the local semantic heuristic and $\gamma_{SSM}$ is parameterized by the $\Delta_{SSM} = 40\%$ baseline.

\section{Conclusion}

The A Priori Boundary demands that we predict the deviation distribution *before* it is measured. By formalizing Aaronson's "compiler diagnostics" as the fundamental constants governing the subjective updating of an embedded QBist agent, we move from post hoc relabeling to predictive mathematics. The lab must now use these formalized Epistemic Horizons to predict the precise collapse distributions for upcoming multimodal or scaled constraint tests.

\end{document}
